%% ma: L11-B14-0012
%% lop: 11
%% chuong: 4
%% bai: 14
%% dang: 11.14.1
%% loai: TN4
%% muc: NB
%% dap_an: "A"
%% nguon: "tai-lieu-hoc-tap-mon-toan-lop-11-theo-chuong-trinh-giao-duc-pho-thong-2018.pdf (Lớp Toán Thầy Duy, Phạm Lê Duy, THPT Chu Văn An), Chương 4 Bài 5, Phần 1, Câu 7 (tr.526)"
%% kien_thuc: "Lớp 11 Bài 14 (phép chiếu song song)"
%% trang_thai: cho-duyet
%% ly_do: "Dạng mới đề xuất (11.14.x), chờ giáo viên duyệt dạng; tài liệu gốc không có lời giải mẫu, đáp án do kho tự giải và kiểm bằng sympy."
\begin{ex}
Cho hình lăng trụ $ABC.A^{\prime}B^{\prime}C^{\prime}$. $A^{\prime}$ là hình chiếu của $A$ trên $(A^{\prime}B^{\prime}C^{\prime})$ qua phép chiếu song song theo phương là đường thẳng nào?
\choice{\True $CC^{\prime}$.}{$CA^{\prime}$.}{$C^{\prime}A^{\prime}$.}{$B^{\prime}A^{\prime}$.}
\loigiai{\sloppy $A^{\prime}$ là hình chiếu của $A$ nên phương chiếu song song hoặc trùng $AA^{\prime}$. Trong các đường đã cho chỉ có $CC^{\prime}\parallel AA^{\prime}$ (cạnh bên của lăng trụ).}
\end{ex}
